% !TeX root = ../../tesis.tex

\section{Estructura de Salidas: JSON Descriptivo y GeoJSON}
\label{anx:apimetro-salidas}

\subsection*{Salida JSON descriptiva (atributos alfanuméricos)}

\begin{lstlisting}[language=json,
    caption={Fragmento de salida JSON descriptiva para líneas de
             transporte.},
    label=lst:json-descriptivo]
[
  {
    "linea_id": 191,
    "nombre": "Metro Linea 1",
    "num_comercial": "1",
    "sistema": "METRO",
    "anio_inauguracion": 1969,
    "color_en": "Pink",
    "color_esp": "F94F8E",
    "tam_km": 16.66467,
    "existe": true,
    "clasificacion": "existente"
  }
]
\end{lstlisting}

\subsection*{Salida GeoJSON ruteable (geometría con métricas)}

\begin{lstlisting}[language=json,
    caption={Estructura GeoJSON ruteable de una estación en
             \texttt{Apimetro}.},
    label=lst:geojson-estacion]
{
  "type": "FeatureCollection",
  "features": [
    {
      "type": "Feature",
      "geometry": {
        "type": "Point",
        "coordinates": [-99.13402, 19.42932]
      },
      "properties": {
        "nombre": "20 de Noviembre",
        "sistema": "MB",
        "alcaldia_municipio": "Cuauhtémoc",
        "es_cetram": false,
        "jerarquia_transporte": "masivo_mediano",
        "tipo_entidad": "estacion"
      }
    }
  ]
}
\end{lstlisting}

\subsection*{Conversión a GeoPackage para uso en \textsc{qgis}}

Para el urbanista que trabaja en escritorio, convertir la salida GeoJSON a
\textit{GeoPackage} requiere una sola línea en Python:

\begin{lstlisting}[language=Python,
    caption={Conversión de GeoJSON a GeoPackage con \textit{GeoPandas}.},
    label=lst:geopandas-conversion]
import geopandas as gpd
gpd.read_file("estaciones.geojson").to_file("estaciones.gpkg",
                                            driver="GPKG")
\end{lstlisting}

El archivo \texttt{.gpkg} resultante es compatible con \textsc{qgis},
\textsc{arcgis} y cualquier cliente que soporte el estándar \textsc{ogc}
GeoPackage \citep{ogc2014geopackage}.
